\documentclass[10pt]{article}
\usepackage[margin=20mm]{geometry}
\usepackage[T1]{fontenc}
\usepackage{lmodern}
\begin{document}
\section*{Tripling: literal UInt64 tree and repeated branch}
\textbf{Result.} The literal source multiplication/addition/subtraction tree decodes to the reviewed weighted field tripling formula. The composed repeated-support theorem derives nonzero original ordinate, successful scoped Some return, canonical output, nonzero output weights and normalized Mumford validity from original input validity and literal branch guards.

\begin{center}\fbox{Raw source constants5,8,10 retained}\par $\downarrow$\par
\fbox{Prove literal reduction with $c\le10$ and no overflow}\par $\downarrow$\par
\fbox{Canonical expression tree decodes to exact field payload}\par $\downarrow$\par
\fbox{Original-input repeated-support field theorem}\par $\downarrow$\par
\fbox{Some canonical, nonzero-weight, normalized-valid output}\end{center}

\textbf{Arithmetic envelope.} Prime $5\le p<2^{31}$, canonical point coordinates and the three coefficients used by triple. For a literal $c\le10$ multiplying canonical $a$, $ca\le10(p-1)<2^{64}$. Thus the actual UInt64 expression $(c*a)\bmod p$ equals $((c\bmod p)*a)\bmod p$. Ordinary canonical products use the existing $p^2<2^{62}$ bound. Constants8 and10 are explicitly not assumed canonical at p5/p7. The tree retains the literal5 inherited from the derivative expression as well.

\textbf{Source route.} This is the retained odd-prime Field route with n=1 and tables=None. The table backend is excluded: unreduced small-prime literals are not claimed to be valid table indices. The word theorem concerns pure values and scoped dispatch, not counters, compiler execution or historical Cargo linkage. tripleWords retains exact source association, and triple\_literal\_tree proves equality to the separately reduced-constant expression tree before the existing expression-composition theorem is invoked.

\textbf{Branch composition.} Original canonical valid point/divisor inputs, four nonzero weights, literal outer K=0, cancellation false and equal-other support guards derive positive coincidence and nonzero original ordinate. The reviewed field theorem supplies output validity and nonzero weights; decode correspondence transports them to words. The constant curve coefficient a0 appears only in original validity and is not read by triple. Outside the repeated-positive branch the scoped wrapper returns None as a restriction marker, not as the production behavior for other branches.

\textbf{Validation.} Main attempt2 and separate axiom readback pass. Literal arithmetic controls cover p5, p7 and the upper-bound representative p2147483647; negative checks establish the constants are indeed noncanonical at the small primes. A full original-input p5 repeated-support fixture invokes the composed theorem; a nonunit p7 fixture invokes complete payload decode. Boundary controls cover cancellation and distinct-other exclusion. Native receipts retain all outcomes; independent source review follows sealing.

The control-development attempts retain theorem-argument inference and finite-field polynomial-coefficient simplification failures. No scientific trials or timing comparisons were run. Cost ratio is null. Next work can use this component in broader dispatch coverage while retaining the existing class-sum and compiler/counter boundaries.
\end{document}
